\documentclass[11pt]{article}
\usepackage[margin=1in]{geometry}
\usepackage{amsmath,amssymb,amsthm,mathtools}
\usepackage[T1]{fontenc}
\usepackage[hidelinks,hypertexnames=false]{hyperref}
\usepackage{booktabs}
\usepackage{enumitem}
\urlstyle{same}
\setlength{\emergencystretch}{3em}
\sloppy

\newtheorem{theorem}{Theorem}
\newtheorem{proposition}[theorem]{Proposition}
\newtheorem{lemma}[theorem]{Lemma}
\newtheorem{corollary}[theorem]{Corollary}
\newcommand{\Q}{\mathbf Q}
\newcommand{\F}{\mathbf F}
\newcommand{\Z}{\mathbf Z}
\newcommand{\Gal}{\operatorname{Gal}}
\newcommand{\Sel}{\operatorname{Sel}}
\newcommand{\Sha}{\mathop{\mathrm{Sha}}}
\newcommand{\ord}{\operatorname{ord}}
\newcommand{\artifact}[1]{\href{run:../#1}{\nolinkurl{#1}}}

\title{Triviality of the Tate--Shafarevich Group of the Elliptic Curve 43a1}
\author{A curve-specific verification note}
\date{2 August 2026}

\begin{document}
\maketitle

\begin{abstract}
Let
\[
 E/\Q:\qquad y^2+y=x^3+x^2,
\]
the elliptic curve with Cremona label \(43\mathrm a1\).  We prove
\(\Sha(E/\Q)=0\).  The proof identifies an exactly normalized Heegner point
over \(K=\Q(\sqrt{-7})\), proves that it generates \(E(K)\), verifies every
residual-image hypothesis in Kolyvagin's odd-primary bound, and uses a full
2-descent for the 2-primary part.  All computer calculations used in the
proof are exact-arithmetic or algorithmically certified; scripts, retained
outputs, and replay commands are linked.
The declared trust boundary is the ordinary one for published theorems and
the exact PARI/GP and eclib algorithms named below.  This is a result about
one specified curve.  It is not a proof of the Birch and Swinnerton-Dyer
conjecture, does not imply the conjecture for any family, and carries no claim
of novelty.
\end{abstract}

\section{Scope and trust statement}

The purpose of this note is deliberately narrow.  Its only main assertion is
the following.

\begin{theorem}\label{thm:main}
For the elliptic curve
\[
 E/\Q:y^2+y=x^3+x^2
\]
with Cremona label \(43\mathrm a1\), one has
\[
 \Sha(E/\Q)=0.
\]
\end{theorem}

The theorem is proved under standard theorem/CAS trust.  ``Theorem trust''
means that the cited published results of Kolyvagin, in the precise form
quoted by Cha, and the standard results on semistable residual
representations, Tate curves, rational torsion, and finite subgroups of
\(\operatorname{GL}_2(\F_p)\) are accepted.  ``CAS trust'' means that the
documented exact arithmetic, full 2-descent, modular-form, modular
parametrization, and saturation algorithms in PARI/GP 2.15.4 and eclib/mwrank
are accepted.  No floating-point recognition, rounded BSD quotient, LMFDB
value of \(\Sha\), or unproved database image label is a proof input.

This trust statement is not decorative.  The exact CM calculation invokes
PARI's modular-form and Taniyama-series constructors.  The Selmer upper bounds
invoke PARI's full 2-descent.  The integral Mordell--Weil basis invokes eclib's
descent and saturation.  The retained artifacts make those calls replayable,
but they do not expand the internals of those systems into a kernel-independent
formal proof.

No novelty claim is made.  The curve is small, its expected arithmetic is in
standard databases, and the ingredients are classical.  The point here is an
explicit proof packet with its hypotheses and trust boundary exposed.

\section{Exact arithmetic data}

Put \(P=(0,0)\).  Direct invariant calculation from the displayed equation
gives
\[
 b_2=4,\quad b_4=0,\quad b_6=1,\quad b_8=1,
 \quad c_4=16,
 \quad \Delta=-43,
 \quad j=-\frac{4096}{43}.
\]
Thus the equation is globally minimal, \(N_E=43\), and \(E\) is semistable,
with nonsplit multiplicative reduction of type \(I_1\) at 43 and good
reduction elsewhere.  In particular \(v_{43}(\Delta_{\min})=1\).

Completing the square gives the 2-division polynomial
\[
 f_2(X)=4X^3+4X^2+1.
\]
It has no rational root and has discriminant \(-16\cdot43\).  Hence
\(E(\Q)[2]=0\), and its splitting field has Galois group
\(S_3=\operatorname{GL}_2(\F_2)\).  Exact counts give
\[
 \#E(\F_2)=5,\qquad \#E(\F_3)=6.
\]
Reduction therefore proves that \(E(\Q)_{\rm tors}=0\).  These elementary
calculations are replayed without external packages by
\artifact{millennium-prize/birch-swinnerton-dyer/43a1/verify_43a1_exact.py};
the retained output is
\artifact{millennium-prize/birch-swinnerton-dyer/43a1/output-exact.txt}.

\section{The normalized Heegner point}

Let
\[
 K=\Q(\sqrt{-7}),\qquad
 \tau=\frac{-37+\sqrt{-7}}{86}.
\]
The discriminant \(-7\) is fundamental and has class number one.  Moreover
\[
 37^2+7=4\cdot43\cdot8,
\]
so 43 splits in \(K\), and the maximal-order Heegner hypothesis holds for
\(X_0(43)\).

Fix the strong Weil parametrization
\[
 \pi:X_0(43)\longrightarrow E,
 \qquad \pi(\infty)=O,
\]
with pullback of the Neron differential equal to the normalized newform
differential.  Let \(y_K\) be the trace to \(K\) of the maximal-order CM point
under this parametrization.  This convention fixes the base point, sign, and
trace normalization used in Kolyvagin's theorem.

\begin{proposition}\label{prop:heegner}
With these conventions,
\[
 \pi(\tau)=P=(0,0),\qquad y_K=P.
\]
\end{proposition}

\begin{proof}
The Hilbert class polynomial is \(H_{-7}(T)=T+3375\).  The primitive forms
\([43,37,8]\) and \([1,37,344]\), corresponding to \(\tau\) and
\(43\tau\), both reduce to \([1,1,2]\).  Therefore
\[
 j(\tau)=j(43\tau)=-3375.
\]
Set \(S=j(z)+j(43z)\) and \(R=j(z)j(43z)\).  At the CM point,
\[
 S=-6750,
 \qquad R=11390625.                                      \tag{1}
\]

The exact certificate constructs the Laurent series \(X(q),Y(q)\) of the
optimal parametrization and verifies
\[
 Y^2+Y=X^3+X^2,
 \qquad \frac{qX'(q)}{2Y(q)+1}=f(q).
\]
It also checks modular degree two, Fricke eigenvalue \(+1\), and trivial
rational torsion.  By Manin--Drinfeld \cite{Manin}, the image of the other
rational cusp is torsion, hence zero; thus the parametrization is the
normalized Fricke quotient, without an unresolved negation or torsion
translation.  The functions \(S\)
and \(R\) descend to rational functions on \(E\), belonging respectively to
\(L(43O)\) and \(L(44O)\).  Their coefficients are recovered over \(\Q\) from
formal \(q\)-series and independently checked through orders 43 and 44.  The
Riemann--Roch zero bound turns those finite checks into identities, rather than
numerical approximations.

Writing
\[
 S=A_S(X)+B_S(X)Y,
 \qquad R=A_R(X)+B_R(X)Y,
\]
the exact values needed at \(X=0\) are
\[
 A_S(0)=-6750,
 \quad B_S(0)=-22750,
 \quad A_R(0)=11390625.                                  \tag{2}
\]
The certificate takes resultants with the Weierstrass equation and computes
the monic gcd of three nonzero projected resultants.  That gcd is exactly
\(X\).  Hence every common point satisfying (1) has \(X=0\).  The curve
equation then gives \(Y=0\) or \(-1\), while (2) selects \(Y=0\).  Thus the
CM fiber is the single point \(P\).  Since the class number is one, the trace
defining \(y_K\) has one summand, proving \(y_K=P\).

Every operation just described is exact integer or rational arithmetic after
the standard PARI constructors are called.  The complete executable
certificate is
\artifact{millennium-prize/birch-swinnerton-dyer/43a1/verify_dminus7_cm_exact.gp},
and its retained successful output is
\artifact{millennium-prize/birch-swinnerton-dyer/43a1/output-dminus7-exact.txt}.
\end{proof}

\section{The Mordell--Weil index over the quadratic field}

The quadratic twist by \(-7\) has global minimal model
\[
 E^{(-7)}:[0,-1,1,-16,-106],
 \qquad y^2+y=x^3-x^2-16x-106.
\]
Its discriminant and conductor are \(-43\cdot7^6\) and \(43\cdot7^2\).
Full multiplication-by-two descent gives
\[
 \dim_{\F_2}\Sel^{(2)}(E/\Q)=1,
 \qquad
 \dim_{\F_2}\Sel^{(2)}(E^{(-7)}/\Q)=0.                 \tag{3}
\]
Because both rational 2-torsion groups vanish, (3), together with the point
\(P\), gives ranks one and zero respectively.  The eigenspace decomposition
for \(\Gal(K/\Q)\) therefore gives \(\operatorname{rank}E(K)=1\).

There is no torsion over \(K\).  The prime 2 splits and 3 is inert in \(K\).
The exact counts
\[
 \#E(\F_2)=5,
 \qquad
 \#E(\F_9)=(3+1)^2-a_3^2=12,
 \qquad a_3=-2,
\]
eliminate odd torsion.  The irreducible cubic \(f_2\) remains without a root
over a quadratic extension, so \(E(K)[2]=0\), which eliminates all 2-primary
torsion.

For \(Q\in E(K)\), the point \(Q-\sigma Q\) is anti-invariant and corresponds
to a rational point on \(E^{(-7)}\).  That twist has rank zero, so
\(Q-\sigma Q\) is torsion and hence zero.  Consequently
\[
 E(K)=E(\Q).                                             \tag{4}
\]
Eclib's full 2-descent and saturation at 1000-bit precision returns
\([0:-1:1]=-P\), reports that the basis is already saturated, and certifies
unconditionally that it is the full Mordell--Weil basis.  Combining this with
(4) proves
\begin{equation}
 E(K)_{\rm tors}=0,
 \qquad E(K)=\Z P.                                      \tag{5}
\end{equation}
In particular, Proposition~\ref{prop:heegner} and (5) imply
\[
 [E(K)_{\rm free}:\Z y_K]=1.                            \tag{6}
\]

The twist descent and point counts are replayed by
\artifact{millennium-prize/birch-swinnerton-dyer/43a1/verify_43a1_K.gp},
with output
\artifact{millennium-prize/birch-swinnerton-dyer/43a1/output-K.txt}.
The full saturation transcript is
\artifact{millennium-prize/birch-swinnerton-dyer/43a1/output-mwrank-saturation.txt}.
The detailed certificate and its trust accounting are in
\artifact{millennium-prize/birch-swinnerton-dyer/43a1/K-mordell-weil-certificate.md}.

\section{Residual surjectivity at every prime}

\begin{proposition}\label{prop:image}
For every rational prime \(p\),
\[
 \overline\rho_{E,p}(G_\Q)=\operatorname{GL}_2(\F_p).
\]
\end{proposition}

\begin{proof}
The case \(p=2\) follows from the irreducible 2-division cubic with nonsquare
discriminant.  For odd \(p\), the type \(I_1\) multiplicative reduction at 43
supplies a nonidentity transvection on \(E[p]\).  This includes \(p=43\): the
Tate parameter has valuation one, so its class is not a 43rd power and wild
inertia supplies the required unipotent.  The nonsplit quadratic twist is
unramified and does not alter inertia.

If a semistable elliptic curve over \(\Q\) has reducible \(E[p]\), then it or
its \(p\)-isogenous curve has a rational point of order \(p\).  Mazur's
rational-torsion theorem therefore gives irreducibility for \(p\geq11\).
For \(p=3,5,7\), direct point counts at 2 and 3 give Frobenius discriminants
\(-4\) and \(-8\).  Their reductions are nonsquares at the indicated primes:
\[
 -4\equiv2\pmod3,
 \qquad -8\equiv2\pmod5,
 \qquad -4\equiv3\pmod7.
\]
Thus \(E[p]\) is irreducible for every odd \(p\).

For \(p\geq7\), Dickson's classification implies that an irreducible subgroup
of \(\operatorname{GL}_2(\F_p)\) containing a transvection contains
\(\operatorname{SL}_2(\F_p)\).  The determinant is the surjective cyclotomic
character, so the image is all of \(\operatorname{GL}_2(\F_p)\).  For
\(p=3,5\), exhaustive exact enumeration handles the exceptional small groups:
after fixing the transvection, every matrix with the required Frobenius trace
and determinant generates a group of order 48 and 480 respectively, exactly
the orders of \(\operatorname{GL}_2(\F_3)\) and
\(\operatorname{GL}_2(\F_5)\).

The point counts, cubic discriminant, and finite matrix enumerations are
replayed by
\artifact{millennium-prize/birch-swinnerton-dyer/verify_cycle261_43a1_residual.py}.
The theorem-level infinite-prime reduction and the characteristic-43 inertia
argument are written out in
\artifact{millennium-prize/birch-swinnerton-dyer/cycle-261-43a1-all-prime-residual-surjectivity.md}.
No Galois-image database is used.
\end{proof}

\section{Odd-primary vanishing}

We use Kolyvagin's inequality exactly as printed in Cha
\cite[Theorem 3, p.~155]{Cha}.  In its standing modular-Heegner setting, if
\(p\) is odd, \(y_K\) has infinite order, and
\(\Gal(\Q(E[p])/\Q)=\operatorname{GL}_2(\F_p)\), then
\begin{equation}
 \ord_p\#\Sha(E/K)\leq 2m_p,                            \tag{7}
\end{equation}
where \(m_p\) is the largest integer for which
\(y_K\in p^{m_p}E(K)\) modulo \(p\)-torsion points, in Cha's wording.  The
Kolyvagin theorem package also gives the finiteness needed for the order in
(7).

The exact wording matters.  Theorem 3 has no condition \(p\nmid D_K\), no
condition \(p\nmid N_E\), and no good-reduction condition at \(p\).  Therefore
it covers both \(p=7\), which ramifies in \(K\), and \(p=43\), where \(E\)
has bad multiplicative reduction.  Those extra local conditions occur in
Cha's later irreducibility-only variant, Theorem 21, not in Theorem 3.  The
published-text audit, including page anchors and stable DOI links, is
\artifact{millennium-prize/birch-swinnerton-dyer/cycle-262-43a1-kolyvagin-theorem3-literature-audit.md}.

By (5), (6), and Proposition~\ref{prop:heegner}, \(y_K=P\) is non-torsion and
\(m_p=0\) for every prime.  Proposition~\ref{prop:image} supplies the full
residual-image hypothesis for every odd \(p\).  Equation (7) therefore gives
\[
 \Sha(E/K)[p^\infty]=0
 \qquad\text{for every odd prime }p.                    \tag{8}
\]

Restriction preserves local triviality and defines
\[
 \operatorname{res}:\Sha(E/\Q)[p^\infty]
 \longrightarrow\Sha(E/K)[p^\infty].
\]
Corestriction after restriction is multiplication by two.  Hence restriction
is injective on odd-primary groups, and (8) yields
\begin{equation}
 \Sha(E/\Q)[p^\infty]=0
 \qquad\text{for every odd prime }p.                    \tag{9}
\end{equation}
The same restriction argument, together with finite generation of \(E(K)\),
also transfers the qualitative finiteness conclusion to \(\Sha(E/\Q)\): its
restriction kernel lies in the finite group \(H^1(\Gal(K/\Q),E(K))\), and its
image lies in finite \(\Sha(E/K)\).

\section{The 2-primary part}

The full multiplication-by-two descent over \(\Q\) returns one basis element
for the everywhere locally soluble 2-covers.  Its explicit affine quartic is
\[
 C:\quad V^2=U^4-2U^2+4U+1.                            \tag{10}
\]
The exact cover map satisfies the Weierstrass equation identically modulo
(10).  The point \((U,V)=(0,1)\) maps to \((1,1)=-3P\), which represents the
nonzero Kummer class of \(P\).  Therefore
\[
 \Sel^{(2)}(E/\Q)=\langle\delta(P)\rangle\simeq\F_2.   \tag{11}
\]
The Kummer sequence is
\[
 0\longrightarrow E(\Q)/2E(\Q)
 \longrightarrow\Sel^{(2)}(E/\Q)
 \longrightarrow\Sha(E/\Q)[2]
 \longrightarrow0.
\]
Since \(E(\Q)=\Z P\), both the first and middle terms have dimension one.
Thus
\begin{equation}
 \Sha(E/\Q)[2]=0.                                      \tag{12}
\end{equation}

The exact GP checker verifies the complete basis size, literal quartic, cover
map, rational point, Kummer class, and 2-division polynomial:
\artifact{millennium-prize/birch-swinnerton-dyer/43a1/verify_43a1_2descent.gp}.
The certificate explanation is
\artifact{millennium-prize/birch-swinnerton-dyer/43a1/2descent-certificate.md}.
The upper-bound assertion in (11) is precisely where trust in PARI's
documented full 2-descent implementation enters.

The group \(\Sha(E/\Q)\) is finite by the preceding section.  Every nonzero
finite 2-primary group contains a nonzero element of order two.  Equation (12)
therefore implies
\begin{equation}
 \Sha(E/\Q)[2^\infty]=0.                               \tag{13}
\end{equation}

\section{Proof of the theorem}

Equation (9) eliminates every odd-primary component of the finite group
\(\Sha(E/\Q)\), and (13) eliminates its 2-primary component.  A finite abelian
group is the direct sum of its primary components.  Hence
\[
 \Sha(E/\Q)=0,
\]
which proves Theorem~\ref{thm:main}.

\section{Reproduction ledger}

Run the following from the repository root.  The commands perform exact
checks; they do not query LMFDB for the claimed conclusion.

\begin{verbatim}
python3 millennium-prize/birch-swinnerton-dyer/43a1/verify_43a1_exact.py
gp -fq millennium-prize/birch-swinnerton-dyer/43a1/verify_43a1_2descent.gp
gp -fq millennium-prize/birch-swinnerton-dyer/43a1/verify_dminus7_cm_exact.gp
gp -fq millennium-prize/birch-swinnerton-dyer/43a1/verify_43a1_K.gp
printf '0 1 1 0 0\n' | mwrank -v 2 -p 1000
python3 millennium-prize/birch-swinnerton-dyer/verify_cycle261_43a1_residual.py
bash scripts/build-paper.sh sha-43a1
\end{verbatim}

The decisive expected terminal markers are
\texttt{ALL\_EXACT\_2DESCENT\_ASSERTIONS\_PASSED=1},
\texttt{FORMAL\_Q\_CERTIFICATE\_PASSED=1},
\texttt{ALL\_EXACT\_K\_RANK\_TORSION\_ASSERTIONS\_PASSED=1}, and
\texttt{ALL\_PRIME\_RESIDUAL\_SURJECTIVITY\_CERTIFICATE=PASS}; mwrank must
state that the points were already saturated and that the full Mordell--Weil
basis was determined unconditionally.

For convenience, the packet index is
\artifact{millennium-prize/birch-swinnerton-dyer/43a1/README.md}.  Stable
external links for the two theorem sources are Cha's article
\href{https://doi.org/10.1016/j.jnt.2004.08.009}{doi:10.1016/j.jnt.2004.08.009}
and Kolyvagin's chapter
\href{https://doi.org/10.1007/BFb0086267}{doi:10.1007/BFb0086267}.
PARI/GP documentation is at
\url{https://pari.math.u-bordeaux.fr/dochtml/html-stable/}, and eclib source
and documentation are at \url{https://github.com/JohnCremona/eclib}.

\section{What has not been proved}

Theorem~\ref{thm:main} is curve-specific.  It does not prove BSD for
\(43\mathrm a1\): no equality of the leading Taylor coefficient with the BSD
arithmetic expression is asserted here.  It does not prove BSD for all
elliptic curves, for all rank-one curves, or for any nontrivial family.  It
does not constitute a Clay Millennium Prize solution or a reduction of that
problem.  Database entries for analytic rank, regulator, Tamagawa product,
and analytic order of \(\Sha\) are corroboration only and are absent from the
logical chain.  No publication-priority or novelty claim is made.

\begin{thebibliography}{9}

\bibitem{Cha}
B. Cha,
\emph{Vanishing of some cohomology groups and bounds for the
Shafarevich--Tate groups of elliptic curves},
J. Number Theory \textbf{111} (2005), 154--178.
\href{https://doi.org/10.1016/j.jnt.2004.08.009}
{doi:10.1016/j.jnt.2004.08.009}.

\bibitem{Kolyvagin}
V. A. Kolyvagin,
\emph{On the structure of Shafarevich--Tate groups},
in Algebraic Geometry (Chicago, 1989), Lecture Notes in Math. 1479,
Springer, 1991, 94--121.
\href{https://doi.org/10.1007/BFb0086267}{doi:10.1007/BFb0086267}.

\bibitem{MazurTorsion}
B. Mazur,
\emph{Modular curves and the Eisenstein ideal},
Publ. Math. IHES \textbf{47} (1977), 33--186.

\bibitem{Serre}
J.-P. Serre,
\emph{Proprietes galoisiennes des points d'ordre fini des courbes
elliptiques},
Invent. Math. \textbf{15} (1972), 259--331.

\bibitem{Silverman}
J. H. Silverman,
\emph{Advanced Topics in the Arithmetic of Elliptic Curves},
Graduate Texts in Mathematics 151, Springer, 1994, Chapter V.

\bibitem{Manin}
Yu. I. Manin,
\emph{Parabolic points and zeta functions of modular curves},
Math. USSR-Izv. \textbf{6} (1972), 19--64.

\end{thebibliography}

\end{document}
